\RequirePackage{iftex}
\ifPDFTeX
\documentclass[a4paper,12pt]{article}
\usepackage[margin=24mm]{geometry}
\usepackage[T1]{fontenc}
\usepackage[utf8]{inputenc}
\usepackage{graphicx}
\usepackage{CJKutf8}
\else
\documentclass[a4paper,12pt]{jsarticle}
\usepackage[dvipdfmx,margin=24mm]{geometry}
\usepackage[dvipdfmx]{graphicx}
\fi
\usepackage{amsmath,amssymb}
\usepackage{enumitem}

\setlength{\parindent}{0pt}
\setlength{\parskip}{0.35em}
\setlist[enumerate]{leftmargin=3em,itemsep=0.25em,topsep=0.35em}

\newcommand{\tenpoints}{\hfill （10点）}
\newcommand{\fivepoints}{\hfill （5点）}
\newcommand{\subtenpoints}{\hfill （10点）}
\newcommand{\blankbox}[2]{\fbox{\rule{0pt}{2.4ex}\makebox[#1]{#2}}}
\newcommand{\problem}[2]{\par\vspace{0.9em}\noindent{\large 問題 #1}\quad（#2点）\quad\ignorespaces}
\newcommand{\answer}{\par\noindent 解答\quad}

\begin{document}

\ifPDFTeX
\begin{CJK}{UTF8}{min}
\fi

\begin{center}
{\Large 数学1A レポート1}　満点 100点
\end{center}

計算問題では途中式も書くこと。証明問題では，どの定理を使ったか書くこと。

\problem{1}{$5\times 2=10$}
次の文を，$\mathbb{R},\mathbb{N},\forall,\exists$ を用いて記号で書き直せ。

\begin{enumerate}[label=(\alph*)]
\item 任意の実数 $x$ に対して，$x$ より大きい自然数 $n$ が存在する。 
\item ある実数 $a$ が存在して， $a^2 = 2 $ となる。 
\end{enumerate}

\answer
\begin{enumerate}[label=(\alph*)]
\item $\forall x\in\mathbb{R},\ \exists n\in\mathbb{N}\ (x<n).$
\item $\exists a\in\mathbb{R}\ (a^2=2).$
\end{enumerate}

\problem{2}{10}
$f(x)=x^2$ とする。このとき
\(
\lim_{x\to 1} f(x)=1
\) 
であることを，次の四角を埋めることで示せ:  $\varepsilon>0$ を任意にとる。このとき, 
\[
\delta=\blankbox{10em}{A}
\]
とおく。$0<|x-1|<\delta$ のとき，
\[
|f(x)-1|
=|x^2-1|
=|x-1||x+1|
<\blankbox{8em}{B}
<\varepsilon.
\]

\answer
$\varepsilon>0$ を任意にとる。
\[
\delta=\min\left\{1,\frac{\varepsilon}{4}\right\}
\]
とおく。$0<|x-1|<\delta$ ならば $|x-1|<1$ であるから，
$0<x<2$ となり，$|x+1|<3$ である。したがって
\[
|f(x)-1|
=|x^2-1|
=|x-1||x+1|
<3|x-1|
<3\delta
\leq \frac{3\varepsilon}{4}
<\varepsilon.
\]
よって，A は $\displaystyle \min\left\{1,\frac{\varepsilon}{4}\right\}$，
B は $3|x-1|$ とすればよい。

\problem{3}{10}
赤道上の各地点の気温を角度 $\theta$ の関数 $T(\theta)$ で表す。$T$ は連続であり，$T(\theta+2\pi)=T(\theta)$ を満たすと仮定する。 このとき，赤道上には互いに反対側にある2地点で，気温が等しいものが存在することを中間値の定理を用いて示せ。

\answer
\[
F(\theta)=T(\theta+\pi)-T(\theta)
\]
とおく。$T$ は連続なので，$F$ も連続である。また $T(\theta+2\pi)=T(\theta)$ より
\[
F(\theta+\pi)=T(\theta+2\pi)-T(\theta+\pi)
=T(\theta)-T(\theta+\pi)
=-F(\theta)
\]
である。

$F(0)=0$ ならば $\theta=0$ と $\theta=\pi$ で気温が等しい。
$F(0)\neq 0$ のとき，$F(\pi)=-F(0)$ であるから，$F(0)$ と $F(\pi)$ は符号が反対である。
中間値の定理より，ある $c\in(0,\pi)$ が存在して $F(c)=0$ となる。
したがって
\[
T(c+\pi)=T(c)
\]
であり，互いに反対側にある2地点で気温が等しい。

\problem{4}{$5\times 2=10$}
次の関数の微分を求めよ。

(1) $\displaystyle  (x^2+1)\sin x$ 
(2) $\displaystyle \frac{e^x}{1+x^2}$ 

\answer
\[
\frac{d}{dx}\{(x^2+1)\sin x\}
=2x\sin x+(x^2+1)\cos x.
\]
\[
\frac{d}{dx}\left(\frac{e^x}{1+x^2}\right)
=\frac{e^x(1+x^2)-2xe^x}{(1+x^2)^2}
=\frac{e^x(x^2-2x+1)}{(1+x^2)^2}.
\]

\problem{5}{10}
$0<a<b$ とする。平均値の定理を用いて，次を示せ。

\[
\frac{\sin b-\sin a}{b-a}\leq 1
\]

\answer
$\sin x$ は $[a,b]$ で連続，$(a,b)$ で微分可能である。
平均値の定理より，ある $c\in(a,b)$ が存在して
\[
\frac{\sin b-\sin a}{b-a}=\cos c
\]
となる。任意の実数 $c$ に対して $\cos c\leq 1$ であるから，
\[
\frac{\sin b-\sin a}{b-a}\leq 1
\]
が成り立つ。

\problem{6}{$5\times 2=10$}
ロピタルの定理を用いて，次の極限を求めよ。

\begin{enumerate}[label=(\alph*)]
\item $\displaystyle \lim_{x\to 0}\frac{\sin x}{x}$ 
\item $\displaystyle \lim_{x\to 0}\frac{e^x-1-x}{x^2}$ 
\end{enumerate}

\answer
\begin{enumerate}[label=(\alph*)]
\item 分子と分母はともに $0$ に近づくので，ロピタルの定理より
\[
\lim_{x\to 0}\frac{\sin x}{x}
=\lim_{x\to 0}\frac{\cos x}{1}
=1.
\]
\item 分子と分母はともに $0$ に近づくので，ロピタルの定理より
\[
\lim_{x\to 0}\frac{e^x-1-x}{x^2}
=\lim_{x\to 0}\frac{e^x-1}{2x}.
\]
右辺も $0/0$ 型なので，もう一度ロピタルの定理を用いると
\[
\lim_{x\to 0}\frac{e^x-1}{2x}
=\lim_{x\to 0}\frac{e^x}{2}
=\frac12.
\]
\end{enumerate}

\problem{7}{$5\times 2=10$}
次の関数の高階微分を求めよ。

\begin{enumerate}[label=(\alph*)]
\item $f(x)=x^2e^x$ の $f^{(3)}(x)$ 
\item $g(x)=\sin x$ の $g^{(2026)}(x)$ 
\end{enumerate}

\answer
\begin{enumerate}[label=(\alph*)]
\item
\[
f'(x)=(x^2+2x)e^x,\quad
f''(x)=(x^2+4x+2)e^x,\quad
f^{(3)}(x)=(x^2+6x+6)e^x.
\]
\item $\sin x$ の導関数は
\[
\sin x,\ \cos x,\ -\sin x,\ -\cos x
\]
を周期 $4$ で繰り返す。$2026\equiv 2\pmod 4$ であるから
\[
g^{(2026)}(x)=-\sin x.
\]
\end{enumerate}

\problem{8}{$10\times 2=20$}
次の関数の マクローリン展開を求めよ。

(1) $\sin x$  (2) $\log(1+x)$

\answer
\[
\sin x
=x-\frac{x^3}{3!}+\frac{x^5}{5!}-\frac{x^7}{7!}+\cdots
=\sum_{n=0}^{\infty}(-1)^n\frac{x^{2n+1}}{(2n+1)!}.
\]
\[
\log(1+x)
=x-\frac{x^2}{2}+\frac{x^3}{3}-\frac{x^4}{4}+\cdots
=\sum_{n=1}^{\infty}(-1)^{n+1}\frac{x^n}{n}
\quad (|x|<1).
\]

\problem{9}{10}
テイラー展開の歴史について調べ，最低５行以上でまとめよ。

\answer
テイラー展開は，関数をある点のまわりで多項式の和として表す考え方である。
この考え方は微分積分学の発展と深く関係しており，ニュートンやライプニッツの時代から研究されていた。
イギリスの数学者ブルック・テイラーは，1715年の著作で現在テイラーの定理と呼ばれる結果を発表した。
その後，マクローリンは $0$ のまわりのテイラー展開を体系的に用いたため，これはマクローリン展開と呼ばれるようになった。
テイラー展開は，三角関数，指数関数，対数関数などを多項式で近似するために重要である。
現在では，解析学だけでなく，物理学，工学，数値計算など多くの分野で用いられている。

\ifPDFTeX
\end{CJK}
\fi

\end{document}
